% ============================================================================
% Fichier  : partie0/P00_vingt_quatre_juin.tex
% Titre    : Le 24 Juin 2026 (prologue opérationnel)
% Auteur   : MINKA MI NGUIDJOÏ Thierry Emmanuel
% Version  : 1.0 -- novembre 2026
% Corrigés : corriges/Corr_P00.tex
% ============================================================================
% Sources : loi n° 2024/017, art. 2, 5, 19, 22, 29, 37-39, 62, 73 (texte
% officiel, scan certifié conforme) ; FAV pour la fuite du 14/03/2025.
% Faits posés (bible) : lettre de M. TAGNE du 24/06/2026 ; demande de l'OPJ du
% 20/03/2025 de différer l'information des clients ; mandat d'audit du
% 30/06/2026 à KORO AUDIT & CONFORMITÉ.
% ============================================================================

\chapter{Le 24 Juin 2026}
\label{chap:prologue}
\chaptermark{Le 24 juin 2026}

\epigraph{%
    « Être en règle, c'est pouvoir le montrer.\\
    Le dire ne suffit pas. »%
}{Principe du cours}

\niveauChapitre{Introductif}{%
    Aucun, sinon la lecture de l'avant-propos
}

\begin{boxobjectifs}
\begin{itemize}
    \item Reconnaître, dans une situation réelle, les questions qui relèvent
          de la protection des données à caractère personnel.
    \item Situer une organisation par rapport à l'échéance du 23~juin~2026
          fixée par l'article~73 de la loi n°~2024/017.
    \item Distinguer une déclaration de conformité d'une conformité
          démontrée.
    \item Coter une exigence \OUI{}, \NON{} ou \NEK{} et justifier la cote.
    \item Situer les métiers qui interviennent, de la réclamation au
          rapport d'audit.
\end{itemize}
\end{boxobjectifs}

\bigskip

% ============================================================================
\section*{Avant de commencer : une lettre, un mercredi matin}
% ============================================================================

Bafoussam, mercredi 24~juin~2026, 8~h~15. Le courrier de NYANGA
CRÉDIT\footnote{NYANGA CRÉDIT, ses dirigeants, ses clients, ses prestataires
et le cabinet KORO AUDIT \& CONFORMITÉ sont fictifs, comme les dates, les
adresses et les montants des cas de ce cours. Les mécanismes et les textes
cités sont exacts.}, établissement de microfinance de vingt-deux agences et
de 210\,000 clients, contient une lettre recommandée de M.~TAGNE Paul,
commerçant au marché A, client depuis 2016.

M.~TAGNE a assisté, en février, à l'audience du tribunal de grande instance
de la Mifi qui jugeait l'affaire de la nuit du 14~mars~2025. Il y a appris
que son dossier figurait parmi les 18\,412 fichiers du dossier
\texttt{Clients} copiés cette nuit-là vers l'étranger. Personne ne l'en
avait informé. Sa lettre demande trois choses~: savoir quelles données
l'établissement détient sur lui et à qui il les a communiquées, obtenir
l'effacement de celles qui ne sont plus nécessaires, et comprendre pourquoi
il n'a pas été prévenu. Elle annonce qu'à défaut de réponse il saisira le
juge des référés.

À 10~heures, le directeur général, M.~TCHOUA, réunit la RSSI, Mme~NDOUMBE,
le DSI, M.~KAMGA, et la juriste, Me~ETOUNDI. Il pose une seule question~:
« Sommes-nous en règle~? » Me~ETOUNDI rappelle que la veille, le 23~juin,
a expiré le délai de dix-huit mois que la loi accordait aux organisations
pour se mettre en conformité. M.~KAMGA assure que « tout est sécurisé
depuis l'incident ». Mme~NDOUMBE répond qu'elle ne sait pas, et qu'elle ne
peut pas le savoir sans un examen.

\begin{boxalert}
\textbf{Violation typique.} Répondre « nous sommes conformes » sur la foi
d'une politique écrite, d'un pare-feu renforcé ou d'une absence de
plainte.

\smallskip
\textbf{Règle~:} la conformité se démontre traitement par traitement, par
des preuves~: un registre, des autorisations, des contrats, des journaux,
des réponses aux personnes. Ce qui n'a pas été examiné reste inconnu, et se
dit inconnu.
\end{boxalert}

% ============================================================================
\section{Ce que l'on sait le 24 juin}
\label{sec:p00-situation}
% ============================================================================

La réunion permet de réunir quelques faits. Ce sont des
\termecle{constatations}~: chacune se rattache à une pièce que quelqu'un
d'autre pourrait vérifier.

\begin{itemize}
    \item Le 14~mars~2025, 3,2~Go du dossier \texttt{Clients} ont été copiés
          vers une adresse étrangère~; l'enquête technique l'a établi, et le
          jugement de février~2026 le reprend.
    \item Le 20~mars~2025, l'officier de police judiciaire chargé de
          l'enquête a demandé par écrit à l'établissement de ne pas informer
          les clients, pour ne pas compromettre l'enquête. Aucun client n'a
          été informé depuis.
    \item L'établissement n'a jamais adressé de demande d'autorisation de
          traitement à l'Autorité prévue par la loi.
    \item Me~ETOUNDI n'a trouvé aucune liste des traitements de
          l'établissement. Le DSI tient un inventaire des applications, qui
          n'est pas la même chose.
\end{itemize}

Les autres phrases entendues pendant la réunion sont des opinions~: « tout
est sécurisé », « les clients ont signé les conditions générales », « le
prestataire de recouvrement est sérieux ». Elles ne sont ni vraies ni
fausses à ce stade~; elles attendent leur preuve.

\begin{boxdef}
\termecle{Donnée à caractère personnel}~: selon l'article~5 de la loi,
« information se rapportant à une personne permettant de l'identifier
directement ou indirectement, notamment par référence à toute forme
d'identifiant ou à un ou plusieurs éléments, propres à son identité
physique, psychologique, génétique, psychique, culturelle,
socio-professionnelle ou économique ». La loi cite parmi les exemples un
nom, une photo, une adresse, un numéro de téléphone, un matricule interne,
une adresse IP et un enregistrement vocal.

\smallskip
\termecle{Traitement}~: toute opération appliquée à de telles données,
« notamment la collecte, l'enregistrement, l'organisation, la conservation
[\ldots], la consultation, l'utilisation, la communication [\ldots], ainsi
que le verrouillage, l'effacement ou la destruction » (article~5).
\end{boxdef}

% ============================================================================
\section{L'échéance du 23 juin}
\label{sec:p00-echeance}
% ============================================================================

La loi n°~2024/017 a été promulguée le 23~décembre~2024~\cite{loi2024017}.
Elle a laissé aux organisations un délai pour adapter leurs pratiques.

\begin{boxjuris}
\textbf{\loicam{n°~2024/017}, article~73}~\cite{loi2024017}~:
\textit{« Les personnes physiques ou morales en charge du traitement des
données à caractère personnel disposent d'un délai de dix-huit (18) mois, à
compter de la date de promulgation de la présente loi, pour se conformer aux
dispositions de celle-ci. »}

\smallskip
\textbf{Portée.} Le délai court du 23~décembre~2024 au 23~juin~2026. À
partir du 24~juin, une organisation qui traite des données ne peut plus
invoquer la période d'adaptation. L'article ne dit pas si les traitements
commencés avant la loi doivent être régularisés rétroactivement~; la
question est discutée au Chapitre~\ref{chap:recherche}.
\end{boxjuris}

La loi s'applique largement. Son article~2 vise tout traitement effectué
par l'État, les collectivités territoriales décentralisées ou toute
personne physique ou morale, tout traitement de données d'une personne
établie, résidant ou en transit au Cameroun, et tout traitement effectué
par un responsable ou un sous-traitant établi au Cameroun. NYANGA CRÉDIT
entre dans chacun de ces cas.

\begin{boxinfo}
\textbf{Ce que la loi ne dit pas.} Elle ne crée pas de délégué à la
protection des données~; elle ne fixe pas de délai chiffré pour informer
d'une violation (l'article~22 dit « sans délai »)~; et de nombreuses
modalités sont renvoyées à des textes d'application. Elle confie enfin
la plus grande part du contrôle à une Autorité de protection des données à
caractère personnel dont la création et le fonctionnement relèvent d'un
décret (article~53, alinéa~2). Le cours lit la loi telle qu'elle est écrite,
et cote \NEK{} ce que ses textes d'application ne permettent pas encore de
trancher.
\end{boxinfo}

% ============================================================================
\section{Trois questions pour un audit}
\label{sec:p00-questions}
% ============================================================================

La question du directeur général, « sommes-nous en règle », est trop large
pour recevoir une réponse. Mme~NDOUMBE la décompose en trois questions,
chacune vérifiable.

\begin{enumerate}
    \item \textbf{Que traitons-nous~?} Quelles données, sur quelles
          personnes, pour quelles finalités, avec quels destinataires et
          pendant combien de temps. C'est l'objet du registre des activités
          de traitement que l'article~29 impose de tenir
          (Chapitre~\ref{chap:cartographie}).
    \item \textbf{Ce que nous faisons est-il permis~?} Le consentement
          ou l'une des trois dérogations de l'article~9, l'autorisation de
          l'article~19, l'information de l'article~21, la durée de
          conservation de l'article~13 (Chapitres~\ref{chap:donnees}
          et~\ref{chap:principes}).
    \item \textbf{Pouvons-nous le prouver~?} Pour chaque exigence, une
          pièce qui l'établit, ou l'aveu qu'elle manque
          (Chapitre~\ref{chap:audit}).
\end{enumerate}

Le 30~juin, le conseil d'administration confie la mission au cabinet KORO
AUDIT \& CONFORMITÉ, de Yaoundé. Son auditrice principale, Mme~BELLO
Aïssatou, demande deux mois de terrain. Elle refuse de commencer par la
sécurité informatique, que M.~KAMGA voulait lui montrer en premier~: « Nous
commencerons par savoir ce que vous traitez. Protéger des données que l'on
n'a pas recensées, c'est protéger au hasard. »

% ============================================================================
\section{Répondre à M. TAGNE}
\label{sec:p00-reponse}
% ============================================================================

La lettre n'attend pas l'audit. Me~ETOUNDI prépare la réponse avec trois
textes en main.

\begin{boxjuris}
\textbf{\loicam{n°~2024/017}, article~39}~\cite{loi2024017}~:
\textit{« La personne concernée peut demander au responsable du
traitement~: les informations permettant de connaitre et de contester le
traitement~; la confirmation que les données à caractère personnel la
concernant font ou ne font pas l'objet de traitement~; la communication des
données à caractère personnel qui la concerne ainsi que toute information
disponible quant à l'origine de celle-ci~; [\ldots] une copie de ses données
à caractère personnel sous une forme intelligible contre paiement d'une
somme ne pouvant excéder le coût de la reproduction. »}

\smallskip
\textbf{Portée.} M.~TAGNE a droit à la liste de ses données, à leur origine
et à leurs destinataires. Le délai de réponse doit être fixé par voie
réglementaire (article~47)~; en attendant, l'auditrice recommande de
répondre dans un délai raisonnable et de le dater.
\end{boxjuris}

La demande d'effacement est plus délicate. L'article~38 oblige à effacer
« sans délai » les données qui ne sont plus nécessaires, mais « sous réserve
des exigences de conservation » prévues par d'autres textes. Or les données
de M.~TAGNE sont aussi des pièces d'une procédure judiciaire et de la
comptabilité de l'établissement. Me~ETOUNDI ne peut ni tout effacer, ni tout
garder~: elle doit séparer, donnée par donnée, ce qui reste nécessaire et ce
qui ne l'est plus. Cette séparation suppose de savoir où se trouvent les
données de M.~TAGNE, ce que personne ne sait encore.

Reste la troisième question~: pourquoi n'a-t-il pas été prévenu~? L'article~22
impose d'informer « sans délai » l'Autorité et la personne concernée dès
connaissance d'une violation. En mars~2025, l'officier de police judiciaire
a demandé le silence pour protéger l'enquête. Deux obligations se sont
rencontrées~; laquelle l'emportait, et à partir de quand la loi
s'appliquait-elle à cette violation antérieure à l'échéance~? Le
Chapitre~\ref{chap:violations} examine la question pour la pratique, et le
Chapitre~\ref{chap:recherche} la présente comme un objet de recherche en
cours.

% ============================================================================
\section{Premières cotes}
\label{sec:p00-cotes}
% ============================================================================

Mme~BELLO ouvre son premier papier de travail le 6~juillet, premier jour de
terrain. Elle cote huit propositions avec ce qu'elle a vu.

\begin{boxdef}
\termecle{Cote \OUI{}, \NON{}, \NEK{}}~: réponse à une question d'audit
fermée sur une exigence. \OUI{}~: la conformité est établie par une preuve.
\NON{}~: l'écart est établi par une preuve. \NEK{} (\textit{not enough
knowledge})~: les preuves disponibles ne permettent pas de trancher.
L'instrument vient du programme de recherche de
l'auditologie~\cite{minka2026nek}.
\end{boxdef}

\begin{boxlecture}{Les huit premières propositions, 6 juillet 2026}
\begin{tableaucours}{L{7.4cm}cL{3.6cm}}{%
    \tableauHeader{Proposition} &
    \tableauHeader{Cote} &
    \tableauHeader{Pièce}
}
NYANGA CRÉDIT traite des données à caractère personnel & \OUI{} & Fichiers
clients, paie \\
NYANGA CRÉDIT est soumise à la loi n°~2024/017 & \OUI{} & Art.~2 \\
Un registre des activités de traitement existe & \NON{} & Aucun produit,
attestation de la juriste \\
Une autorisation de l'Autorité a été demandée & \NON{} & Attestation de la
direction \\
Les clients ont consenti aux traitements & \NEK{} & Formulaires non encore
examinés \\
Les données des clients sont protégées & \NEK{} & Aucune mesure encore
testée \\
Les prestataires sont liés par un contrat conforme & \NEK{} & Contrats non
encore reçus \\
Des données sont transférées hors du Cameroun & \NEK{} & Hébergement du
logiciel bancaire non encore examiné \\
\end{tableaucours}

\medskip
\textbf{Décompte.} Deux \OUI{}, deux \NON{}, quatre \NEK{} sur huit.
\hfill\lecture{0{,}25}{0{,}50}{1{,}50}

\medskip
\textbf{Lecture.} La moitié des propositions est inconnue. Ce n'est pas un
mauvais résultat~: c'est l'état normal d'un audit qui commence. Les deux
\NON{} sont, eux, établis dès le premier jour, et ils suffisent à répondre au
directeur général que l'établissement n'est pas en règle sur ces deux
points.
\end{boxlecture}

Les trois mesures se lisent ensemble. $\mu$ est la part des \OUI{}, $U$
celle des \NEK{}, $H$ l'entropie de la répartition, en bits. Le
Chapitre~\ref{chap:audit} en donne la règle de lecture, et le
Chapitre~\ref{chap:recherche} en discute les limites.

% ============================================================================
\section{Deux autres organisations}
\label{sec:p00-autres}
% ============================================================================

Deux autres organisations accompagnent NYANGA CRÉDIT dans ce cours.

Le \textbf{CENTRE MÉDICAL MBOLO}, à Bertoua, établissement privé de
120~lits, traite des données de santé, qui comptent parmi les données
sensibles de l'article~5. Son dossier patient informatisé est édité et
hébergé par un prestataire de Douala~; il prépare un service de
téléconsultation pour les villages de la Kadey, et il a découvert en mai
qu'un dossier de patiente avait été consulté par des agents sans lien avec
ses soins. Sa directrice, le Dr~ABENA, a commandé un audit de préparation
avant l'échéance.

\textbf{NOUN PAY}, à Bafoussam, fintech de paiement mobile et de
microcrédit, héberge tout son système d'information chez un fournisseur de
nuage à l'étranger et décide de l'octroi de ses crédits par une notation
automatique. Un client conteste un refus. Les transferts de données hors du
Cameroun et les décisions automatisées sont au centre de son audit.

% ============================================================================
\section{Les métiers de l'affaire}
\label{sec:p00-metiers}
% ============================================================================

\begin{boxmetier}{Auditeur de la protection des données}
Mme~BELLO examine, prouve et conclut. Elle ne corrige pas elle-même~: elle
dit ce qui est établi, ce qui manque, et ce qu'il faut faire. Son
indépendance à l'égard de la direction qu'elle audite conditionne la valeur
de ses conclusions.
\end{boxmetier}

\begin{boxmetier}{Délégué à la protection des données}
La loi n'impose pas cette fonction. Beaucoup d'organisations la créent
pourtant, sur le modèle européen, pour qu'une personne tienne le registre,
conseille les équipes et réponde aux personnes. NYANGA CRÉDIT n'en a pas~:
la lettre de M.~TAGNE est arrivée sur le bureau du directeur général.
\end{boxmetier}

\begin{boxmetier}{Juriste conformité}
Me~ETOUNDI lit les textes, répond aux personnes et mesure le risque de
contentieux. Elle doit savoir dire où la loi est claire, où elle renvoie à
des textes non publiés, et où deux obligations se contredisent.
\end{boxmetier}

\begin{boxmetier}{RSSI}
Mme~NDOUMBE tient la sécurité des données, exigée par les articles~22, 24
et~27. Elle sait aussi que la sécurité n'est qu'une partie de la
conformité~: des données bien protégées peuvent être collectées sans droit.
\end{boxmetier}

\begin{boxmetier}{Responsable de traitement}
M.~TCHOUA, pour l'établissement, détermine les finalités et les moyens des
traitements. C'est lui que la loi oblige, et lui qui répond des sanctions
des articles~54 à~71.
\end{boxmetier}

% ============================================================================
\section{Cinq réflexes}
\label{sec:p00-reflexes}
% ============================================================================

\begin{enumerate}
    \item \textbf{Lire le texte officiel.} On cite la loi telle qu'elle est
          écrite, pas telle qu'un résumé ou le modèle européen la décrit.
    \item \textbf{Recenser avant de juger.} On ne peut pas dire si un
          traitement est permis tant qu'on ne sait pas qu'il existe.
    \item \textbf{Demander la preuve.} Une déclaration n'établit rien~; une
          pièce établit quelque chose, et l'on dit quoi.
    \item \textbf{Coter l'inconnu.} Ce qui n'a pas été examiné est \NEK{},
          jamais conforme par défaut.
    \item \textbf{Protéger les données de l'audit.} L'auditeur manipule
          lui-même des données personnelles~: il en prend le moins
          possible et les protège comme il l'exige des autres.
\end{enumerate}

\begin{boxresume}
\begin{itemize}
    \item Depuis le 24~juin~2026, le délai de mise en conformité de
          l'article~73 est expiré~: la loi n°~2024/017 s'applique
          pleinement.
    \item La conformité se démontre par des preuves, traitement par
          traitement~; une déclaration ne la démontre pas.
    \item La question « sommes-nous en règle » se décompose en trois~: que
          traitons-nous, est-ce permis, pouvons-nous le prouver.
    \item À NYANGA CRÉDIT, l'absence de registre et de demande
          d'autorisation est établie dès le premier jour~; la moitié des
          propositions reste \NEK{}.
    \item Une violation antérieure à l'échéance, couverte par une demande de
          silence de l'enquête pénale, pose une question d'articulation des
          normes que le cours traite à part.
\end{itemize}
\end{boxresume}

\begin{boxinfo}
\textbf{Points de contrôle.} Vous avez compris ce prologue si vous savez~:
\begin{itemize}
    \item dater l'échéance de l'article~73 et dire ce qu'elle change~;
    \item séparer, dans un compte rendu de réunion, les constatations des
          opinions~;
    \item coter une exigence \OUI{}, \NON{} ou \NEK{} et citer la pièce qui
          fonde la cote.
\end{itemize}
\end{boxinfo}

\begin{boxexo}[title={\ding{45}\quad Exercice~0.1. Donnée personnelle ou non \niveaubase}]
Dites si chaque élément est une donnée à caractère personnel au sens de
l'article~5, et pourquoi~: (a) le numéro de compte de M.~TAGNE~; (b) le
nombre total de clients de l'agence de Dschang~; (c) l'adresse IP d'où un
client se connecte à l'application mobile~; (d) l'enregistrement d'un appel
au service client~; (e) le taux d'intérêt moyen des crédits de 2025.
\end{boxexo}

\begin{boxexo}[title={\ding{45}\quad Exercice~0.2. Constatation ou opinion \niveaumoyen}]
Classez chaque phrase de la réunion du 24~juin comme constatation ou
opinion, et proposez une cote \OUI{}, \NON{} ou \NEK{} avec la pièce qui la
fonderait~: (a) « tout est sécurisé depuis l'incident »~; (b) « les clients
ont signé les conditions générales »~; (c) « aucune demande d'autorisation
n'a été envoyée »~; (d) « le prestataire de recouvrement est sérieux ».
\end{boxexo}

\begin{boxexo}[title={\ding{45}\quad Exercice~0.3. Répondre à M. TAGNE \niveaumoyen}]
\begin{enumerate}[(a)]
    \item Rédigez la réponse à la demande d'accès de M.~TAGNE, en vous
          appuyant sur l'article~39.
    \item Quelles données pouvez-vous effacer, et lesquelles devez-vous
          garder~? Justifiez avec l'article~38.
    \item Que répondez-vous à sa troisième question sans engager
          imprudemment l'établissement~?
\end{enumerate}
\end{boxexo}

\begin{boxexo}[title={\ding{45}\quad Exercice~0.4. Le silence demandé \niveauexpert}]
Le 20~mars~2025, l'officier de police judiciaire demande à NYANGA CRÉDIT de
ne pas informer ses clients.
\begin{enumerate}[(a)]
    \item Quelles obligations entrent en concurrence~?
    \item La loi n°~2024/017 s'appliquait-elle déjà à cette violation~?
          Distinguez la promulgation et l'échéance de l'article~73.
    \item Quelle trace écrite l'établissement aurait-il dû garder pour se
          défendre aujourd'hui~?
\end{enumerate}
\end{boxexo}

\bigskip

\begin{boxinfo}
\textbf{Pour aller plus loin.}
\begin{itemize}
    \item La loi n°~2024/017~\cite{loi2024017}, dans son texte officiel.
    \item D.~J. Solove, \textit{Understanding Privacy}~\cite{solove2008},
          pour la diversité de ce que l'on appelle la vie privée.
\end{itemize}
\end{boxinfo}

\begin{boxinfo}
\textbf{Transition.} Avant de recenser les traitements de NYANGA CRÉDIT, il
faut savoir ce qu'est exactement une donnée, un traitement, un responsable
et un sous-traitant, et quand une donnée devient sensible. C'est l'objet du
Chapitre~\ref{chap:donnees}.
\end{boxinfo}
